\section{Experimental Setup}

\subsection{Research Questions}

\begin{itemize}
  \item \textbf{RQ1:} Are The Pile's domains measurably distinguishable by cognitive task pattern proxies? (h-m1)
  \item \textbf{RQ2:} Do domain exposure trajectories show sufficient within-family variation? (h-e1)
  \item \textbf{RQ3:} Does domain-specific exposure predict benchmark-specific capability improvement? (h-m2, h-m3)
\end{itemize}

\subsection{Dataset and Model Family}

\textbf{Pre-training corpus:} The Pile \cite{Gao2020Pile} --- 800GB, 22 domains. \textbf{Model family:} Pythia \cite{Biderman2023Pythia} --- 16 models (70M--12B), 154 checkpoints each; focus on 70M, 1B, 6.9B. Benchmarks: MMLU (5-shot), HellaSwag (10-shot), ARC-Challenge (25-shot), WinoGrande (5-shot) via lm-evaluation-harness v0.4.

\subsection{Baselines}

\begin{table}[h]
\centering
\caption{Pre-registered baselines for the panel regression framework.}
\label{tab:baselines}
\begin{tabular}{ll}
\toprule
\textbf{Baseline} & \textbf{Description} \\
\midrule
Scale-only & Panel regression with only $\log(\text{params})$ --- null: domain adds nothing beyond scale \\
Permutation null & Shuffled domain labels (1,000 permutations) --- null $R^2$ distribution \\
Uniform-$\beta$ & Shared domain coefficients across all benchmarks --- null: no specificity \\
\bottomrule
\end{tabular}
\end{table}

\textbf{Note:} The baselines above are defined as part of the pre-registered h-m3 panel regression framework and are implemented in the validated pipeline. Due to the data limitations documented in Section~\ref{sec:results} (Books3 zero-exposure, insufficient evaluation cache), none of these baselines could be executed in the current study. They remain ready to run when the prerequisite data gaps are resolved.

\subsection{Compute}

Benchmark evaluation: 5$\times$ H100 NVL GPUs ($\sim$30 min/checkpoint for 6.9B). Domain lookup: 3.5 min for 600K documents on CPU.
